\section{Water-Table Management and Drainage-System Design}

The frontier in drainage water-quantity management is defined by a single
shift: from passive, one-way conveyance toward active, location-specific
control of the water table and the coupled irrigation--drainage balance. The
enabling ideas---water-table-control structures, automation, forecast-based
operation, and re-evaluated physical design---are largely in place, but each
remains a hypothesis awaiting the long-term, multi-site field and simulation
evidence needed to quantify its yield and water-quality trade-offs. What
follows organises these open questions around active water-table control, the
control systems that would automate it, and the design criteria that would
support it.

\subsection{Controlled drainage, subirrigation, and drainage water management}

Integrating controlled drainage (CD) with irrigation is repeatedly framed less
as settled practice than as a research frontier: the recurring call is to move
from passive, ``unilateral'' drainage toward active control of the water table
and the coupled irrigation--drainage balance
\parencite{Ayars2006b,Kesicka2023b,Ritzema2023}. \textcite{Ayars2006b} argue
that water-table-control structures should be required in new designs and
considered for retrofits, and that new criteria and management methods are
needed---claims that require long-term field studies, possibly combined with
modelling, to validate---echoed by calls to build in control options and expand
changeable CD networks \parencite{Dukhovny2018,Mahmoud2021,Strock2011}. The
quantitative payoff of active management remains open: \textcite{Watkins2024}
hypothesise that active outlet management during extremely wet periods could
minimise wet-stress and raise yields; \textcite{Saadat2017} emphasise the
cost--benefit trade-off; and \textcite{Helmers2012} suggest that CD yield
reductions could be mitigated by greater management while weighing installation
cost, management burden, and applicable acreage. Seasonally differentiated
control---free drainage in wet seasons, CD in dry seasons, with supplemental
irrigation for extremes---is an explicit hypothesis to test
\parencite{Asgari2021}.

Coordinating irrigation with drainage control is a persistent unsolved problem:
irrigation must be reduced when CD is implemented, requiring coordination among
authorities and farmers \parencite{Wahba2005,Ritzema2011,Ritzema2016}, and CD
requires genuinely integrated management with careful site selection
\parencite{Dukhovny2018}. Subirrigation embedded in regional water systems is a
distinct, largely unresolved frontier: regional water needs should be known
before responsible incorporation, with optimised pumping and cautious use of
alternative resources \parencite{deWit2022}, and SWAP-type modelling can
pre-screen sites and forecast consequences before implementation
\parencite{deWit2024}. Seepage effects on CD performance need coordinated
experimental and modelling research, including guard lines to evaluate lateral
seepage \parencite{Youssef2021,Pease2017,Ale2009}, while excess-moisture stress,
spring drawdown, and spring nitrogen loading remain priorities
\parencite{Craft2018,Lahdou2019}. Finally, longer-term studies must disentangle
whether soil-property changes under drainage--subirrigation arise from the
practice itself or from irrigation-water quality \parencite{Kaur2023}, and
complete water and nutrient budgets within and across field boundaries---plus
effects on runoff, preferential flow, and recharge---remain to be tracked
\parencite{Shedekar2021,Ross2016,King2022}.

\subsection{Automation, forecast-based, and data-driven control}

Automation and data-driven control are identified as the next technological
step, but the control systems themselves are still to be built and validated.
\textcite{Youssef2023} call for a new generation of control structures and
decision-support tools, sensor-based systems tailored to shallow-water-table
soils, and systems that adapt to changing weather; \textcite{Ayars2015} extend
this to adaptive-control decision support. Concrete methodological agendas
include dynamic modelling and data assimilation of the field water balance,
feasible formulations of the control problem, and fusion of control automation
with physical infrastructure \parencite{Ikonen2023,Isomäki2024}, automatic
water-depth control for aerobic rice \parencite{Arif2021}, and---anticipated
decades earlier---farm decision aids that jointly incorporate water quantity,
quality, crop growth, and economics \parencite{THOMAS1995}.

Anticipatory, forecast-based operation is the emerging synthesis of these
threads. Forecast-driven management should incorporate evapotranspiration,
subirrigation, and spatially variable seepage, and must explicitly address
forecast uncertainty \parencite{VandeCraats2021}; collector-ditch water levels
and control schedules should be optimised \parencite{Salla2022,Strock2025b};
and control should be paired with local crop and climate characteristics
without overlooking the risk that rapid water-table decrease raises irrigation
demand in dry summers
\parencite{Güngör2023,Dou2023,Golmohammadi2016,Kim2023,Wang2020c}. The open
questions are thus as much about coordinating forecasts, sensors, and structures
into a single adaptive system as about any individual component.

\subsection{Novel structures, design criteria, and drainage intensity}

Novel structures and materials define a first set of open questions: improved
HYDRUS-2D boundary conditions and field verification of siphon drainage
\parencite{Ren2023}, long-term measurement of working time against construction
cost \parencite{Tao2025b}, and new structures that adjust water-holding capacity
while inhibiting rapid vertical and bypass flow \parencite{Li2021}. Multi-level
and bi-level configurations warrant study of actively managed shallow and deep
drains \parencite{Hornbuckle2007,Verma1998,Nelson2011b}, and envelope,
perforation, and screen design remain unresolved across gradients, soils, and
textures \parencite{Alavi2021,Wu2022,Oyarce2017,HajAmor2019}. Optimal layout and
burial-depth determination in saline and problem soils recur in recent studies,
which stress ditching depth, restrictive-layer breaking, amendments, and
construction cost within multi-scheme decision models
\parencite{Tian2023,Tian2024,Tian2023c,Zhang2024d,Zhai2023,Han2023,Qian2025b,Ren2021,Zheng2025,SerranoCoronel2018,Zegeye2021}.
Analytical advances are also sought---formulas for sloped and multilayer soils
\parencite{Tao2016,Tao2019}, partially penetrating and trapezoidal ditches
\parencite{Chahar2010,Afruzi2014}, and mass-transfer and three-dimensional flow
effects \parencite{Zhang2024}---alongside CAD/GIS-based design software and
standardised drainage-coefficient definitions to enable comparison
\parencite{GALLICHAND1992b,Skaggs2017}.

Many site studies conclude that established design criteria themselves require
re-evaluation. New criteria that protect crops while minimising drain-water
salinity and volume are sought \parencite{Christen2001}, along with updated
depths and spacings at specific sites
\parencite{Muhammad2021,Malota2016,Bahçeci2009,Kahlown2002}, review of
recommended spacings paired with controlled-traffic farming
\parencite{Mante2022,Mante2018}, cost-effectiveness of narrow spacings in
claypan soils \parencite{Nelson2017}, and standard depth-and-spacing guidelines
tested in situ before extension \parencite{HajAmor2019}. Design should further
account for changing crop water needs, climate, and irrigation modernisation
\parencite{Rui2024,Akram2013,Ziaeifar2025,Ritzema2023,Ritzema2007,Ritzema2006,Ritzema2011}.

Drainage intensity emerges as a systems question rather than a per-field
specification. Reduced spacing would require increased network capacity and
heighten reliance on best management practices
\parencite{Singh2009,Dialameh2025}; a systems approach should integrate
drainage, crop--soil processes, and nutrient-loss reduction
\parencite{Castellano2019}; and DRAINMOD-based analysis should develop optimum
rates and spacings across locations \parencite{Strock2011}. Which soil
parameters are most critical, and what confidence intervals each requires,
remain fundamental unanswered questions
\parencite{Ayars2015,SHIRMOHAMMADI1995,HARDJOAMIDJOJO1982}. Underlying all of
these is a methodological prerequisite: design guidance must rest on long
records, since drainage design should not rely on short-term simulations
\parencite{Wang2006b}, and long runs with hourly input are needed to assess
wider spacings, interannual variability, and shallow-water-table control
\parencite{Ale2009,Häggblom2019,Malota2015,Salo2019,Kladivko2023,Chang2012}.
